\chapter*{Appendix C: Application 2: Runtime Virtual Keyboard Engine and Live Tracking}
\label{ch:appendix_detector_runtime}

\section*{C.1 Overview of the Runtime Engine}
\label{sec:appendix_detector_overview}
Below are the screenshots of Application~2, which is the runtime virtual keyboard application. This application is responsible for tracking AprilTag markers, tracking hand joint positions using MediaPipe, detecting finger touches using our PyTorch LSTM model, and executing the assigned keyboard shortcuts.

Figure~\ref{fig:appendix_detector_home} represents the opening screen where the user can choose the webcam, select the layout XML file, and pick up an action profile.

\begin{figure}[htbp]
\centering
\includegraphics[width=0.88\textwidth]{figures/detector_home_screen.png}
\caption{Application~2 home screen for camera device selection and layout loading.}
\label{fig:appendix_detector_home}
\end{figure}

\newpage

\section*{C.2 Action Rebinding and Profile Management}
\label{sec:appendix_detector_rebinding}
Figure~\ref{fig:appendix_detector_edit} illustrates the window that allows rebinding of actions. It is possible to assign any desired action to the key in question, either from standard shortcuts (such as Ctrl+C and Alt+Tab) or from terminal scripts. The same XML file is used in both the applications, where key actions and detector settings are saved in the \texttt{<DetectorActions>} and \texttt{<DetectorSettings>} parts of the layout XML. There is no need for additional files anymore, and it is possible to alternate XML action layouts on one printed sheet of paper without printing again.

\begin{figure}[htbp]
\centering
\includegraphics[width=0.92\textwidth]{figures/detector_layout_edit_mode.png}
\caption{Application~2 action rebinding interface for assigning keystrokes and shell commands to macro keys.}
\label{fig:appendix_detector_edit}
\end{figure}

\newpage

\section*{C.3 Live Camera Tracking in Idle and Contact States}
\label{sec:appendix_detector_live_states}
Figure~\ref{fig:appendix_detector_idle} presents the camera live stream without the presence of any hand. The software detects the four AprilTags placed in the corners and determines the homography matrix $H$ to draw the boxes on the camera feed.

\begin{figure}[htbp]
\centering
\includegraphics[width=0.88\textwidth]{figures/detector_playmode_without_hand_just_idle.png}
\caption{Application~2 idle monitoring state showing detected AprilTags, homography alignment, and key bounding boxes.}
\label{fig:appendix_detector_idle}
\end{figure}

Figure~\ref{fig:appendix_detector_touch} demonstrates what occurs when a finger comes into contact with a key. MediaPipe is capable of following the 21 joints of the hand in real time. Once the index finger makes contact with the paper, the PyTorch LSTM model is able to predict a touch probability of 96\% for the index finger, highlights the key in green on the screen, and sends the keystroke to the computer.

\begin{figure}[htbp]
\centering
\includegraphics[width=0.88\textwidth]{figures/detector_playmode_exacty_index_finger_touch_happening_on_a_key.png}
\caption{Application~2 live contact detection showing MediaPipe skeleton, 96\% index touch probability, and highlighted key impact.}
\label{fig:appendix_detector_touch}
\end{figure}

\newpage

\section*{C.4 Production Run Mode and Action Event Log}
\label{sec:appendix_detector_run_mode}
Live Run Mode is illustrated in Figure~\ref{fig:appendix_detector_run}. A line graph displays the touch probability for each finger over time. As seen from the screen, the program runs efficiently at 11.9~FPS with inference time being just 2.3~ms on the CPU. The list on the right displays recent keystrokes.

\begin{figure}[htbp]
\centering
\includegraphics[width=0.92\textwidth]{figures/detector_run_mode_multipl_touches_happend.png}
\caption{Application~2 production run mode showing real-time multi-finger probability history, pipeline FPS, and action event log.}
\label{fig:appendix_detector_run}
\end{figure}

\newpage

\section*{C.5 Runtime Engine Preferences and Detection Parameters}
\label{sec:appendix_detector_settings}
Figure~\ref{fig:appendix_detector_settings_1} and Figure~\ref{fig:appendix_detector_settings_2} illustrate the settings interface. The user can customize the touch sensitivity ($\tau_{\text{touch}}$), delay in between touches, and camera update rates.



\includegraphics[width=0.48\textwidth]{figures/detector_settings_page_1.png}

\includegraphics[width=0.48\textwidth]{figures/detector_settings_page_2.png}

\caption{Touch sensitivity and classification thresholds.}
\label{fig:appendix_detector_settings_1}


